\documentclass[../DoAn.tex]{subfiles}

\begin{document}

\section*{Tóm tắt}

Đồ án nghiên cứu các phương pháp trích xuất, lập chỉ mục, truy xuất và hỏi đáp
thông tin trên tài liệu PDF. Trọng tâm của đồ án không chỉ là tạo câu trả lời, mà
còn là giữ lại cấu trúc và thông tin nguồn của tài liệu PDF để người dùng có thể
kiểm chứng kết quả. Hệ thống được tổ chức thành hai luồng liên kết với nhau.
Luồng tiếp nhận và trích xuất tài liệu phân tích trang PDF, định tuyến các vùng
nội dung, trích xuất văn bản, tái dựng bảng và lưu siêu dữ liệu nguồn. Luồng xử
lý truy vấn xây dựng chỉ mục từ dữ liệu đã trích xuất, truy xuất bằng chứng, kiểm
tra mức độ đầy đủ của bằng chứng, sinh câu trả lời và trả về dẫn chứng nguồn.

Hệ thống kết hợp truy xuất theo từ khóa, truy xuất ngữ nghĩa, truy xuất kết hợp và
hỏi đáp tăng cường bằng truy xuất. Đối với bảng, đồ án đánh giá bộ trích xuất
bảng mặc định, Table Transformer và thiết kế Hybrid TATR kết hợp cấu trúc hình
học của bảng với các hộp từ lấy từ lớp văn bản PDF. Các thực nghiệm được tổ
chức theo nhiều tầng, gồm tiếp nhận và trích xuất PDF, phân tích bố cục, OCR,
tái dựng bảng, truy xuất, chia đoạn có nhận thức bảng, hỏi đáp có căn cứ và các
thí nghiệm loại bỏ thành phần.

Kết quả thực nghiệm trên tập đánh giá QCDT cho thấy quy trình đề xuất có thể
truy xuất bằng chứng liên quan và sinh câu trả lời có dẫn chứng trong phạm vi
các thiết lập đã đánh giá. Tuy nhiên, kết quả trên QASPER cho thấy hệ thống
vẫn còn hạn chế với tài liệu khoa học dài và các câu hỏi cần tổng hợp nhiều bằng
chứng.

\end{document}
